\section{Discussion and Limitations}
\label{sec:discussion}

\subsection{Discussion}
The results support treating explanation, recourse and fairness as parts of one decision process rather than as
separate analyses. An explanation that is the same for every applicant, as whole-model LIME turned out to be, gives the
applicant and the officer nothing to act on; restricting LIME to the applicant's neighbourhood made it agree with SHAP
and made the two explanations shown in the application consistent with each other.

For recourse, the main finding is that validity against the current model is not enough. Three quarters of the
options that the standard search found for rejected applicants failed on at least one model retrained on the same
data. Every safeguard we added (the safety margin, the plausibility check and agreement of retrained models) trades
coverage for durability, and no single setting is best on both. The hybrid resolves this for the applicant: it keeps
the coverage of the standard search, lists the options most likely to last first, and labels each option honestly. In
deployment, the same search models can be refreshed whenever the production model is retrained.

The age experiment illustrates why fairness questions cannot always be settled by removing a feature. Without the
age column, the model reconstructed age from other features less accurately and penalized young applicants more. We
therefore treat the policy questions around age (a minor needs a guardian; age should not otherwise restrict an adult
who earns) as explicit eligibility rules, and keep age in the model as an audited risk factor. Whether the remaining
gap, which follows real differences in repayment, is acceptable is a policy decision for the lender and the regulator;
the system makes the gap, its cause and the cost of closing it visible.

\subsection{Limitations}
\emph{Retraining as a proxy for model change.} Robustness is measured against bootstrap-retrained models on the same
data. This captures the variation of retraining, but not drift in the applicant population over time, which the dataset
cannot show because it has no real application dates. The robustness figures also vary between runs: after we rebuilt
every result with training-only cleaning statistics, the recourse shares moved by 1 to 5 percentage points while every
conclusion held.

\emph{Small DiCE sample.} Because of its run time, DiCE was evaluated on 20 applicants per tier, and its results
should be read as indicative.

\emph{Model-based, not causal, recourse.} The engine assumes that a change affects only the features it modifies, and
the waiting option assumes that the applicant keeps the same job and makes no new credit applications. A causal model
of how these quantities interact~\cite{karimi2022survey} would make the advice more reliable.

\emph{Unequal recourse across groups.} The share of rejected applicants who receive an option differs across age
bands; for the engine with loan levers only it ranges from 50\% (25--35) to 72\% (45--55). Several groups in our
sample are small (19 to 116 applicants), so these differences need to be monitored on larger data before conclusions
are drawn.

\emph{No user study.} Understandability is measured through readability, the amount of information shown and plain
wording, not with applicants or loan officers. An expert review based on usability heuristics has been prepared, and a
study with users is left to future work.

\emph{Illustrative policy and data.} The eligibility thresholds, the guardian rule and the loan pricing are
illustrative and have not been reviewed against lending regulation. A minor's application is assessed on the
guardian's age, income and employment combined with the applicant's other details, whereas a lender would assess the
guardian's full credit record. The data come from one lender and one market, contain no applicants under 20, and may
not represent other populations, including Indian borrowers.

\emph{Order of the experiments.} The age experiment was run before the cleaning statistics were restricted to the
training rows and was not repeated, because its conclusion did not depend on it; its current-model row therefore
differs slightly from the rebuilt audit (0.349 against 0.356).
